\section*{Erd\H{o}s Problem 1199 --- GPT 6 Astra Ultra partial attempt}


\subsection*{FORMAL RESTATEMENT}
For every map $c:\mathbb N\to\{0,1\}$, is there an infinite set $A\subseteq\mathbb N$ and $i\in\{0,1\}$ such that
\[c(a+b)=i\qquad\text{for all }a,b\in A?\]
The requirement includes $a=b$, so all doubled elements $2a$ are included.

\subsection*{CONTEXT AND SCOPE}
This replaces the earlier statement-only record with an elementary mathematical attempt. The arguments below establish only their stated partial results; no novelty or full solution is claimed.

\subsection*{ATTACK PLAN AND WORK}
Infinite Ramsey extraction reduces the difficulty to a precise mismatch.
First choose an infinite set $X$ on which the doubled values $c(2x)$ have
one color $d$, by pigeonholing. Color each unordered pair of distinct members
of $X$ by $c(x+y)$. The nested-tail proof of infinite Ramsey's theorem yields
an infinite $A\subseteq X$ whose distinct-pair sums all have one color $e$.
If $d=e$, this $A$ solves the full problem, including repeated summands.

The mismatch $d\ne e$ cannot simply be ruled out on the extracted set.
For the coloring $c(n)=v_2(n)\bmod2$ and $A=\{2^{2j}:j\ge1\}$, distinct-pair
sums have even 2-adic valuation, while doubled elements have odd valuation.
This example refutes the proposed inference from off-diagonal homogeneity;
it is not a counterexample to the original existence question for that coloring.

The full conclusion does hold for every eventually periodic coloring.
If the period is $q$ beyond $N$, take all sufficiently large positive multiples
of $q$ as $A$. Every sum, including $2a$, lies in the same residue class zero
modulo $q$ beyond $N$, and hence has the same color.

\subsection*{VERIFICATION AND REMAINING GAP}
For an arbitrary coloring, force agreement between the diagonal and
off-diagonal colors or build another set overcoming the mismatch. Ramsey's
theorem on distinct pairs alone is insufficient.

\subsection*{SOURCES}
https://www.erdosproblems.com/1199

\subsection*{FINAL}
\textbf{UNRESOLVED.} The full source problem remains outside the scope of the proved partial results.

\subsection*{COMPLETION ESTIMATE}
$8\%$. This is a conservative subjective estimate toward the whole problem, not a probability that the argument is correct or a claim that a fixed fraction of the problem has been solved.
